%% ════════════════════════════════════════════════════════════════════════════
\section{Extension to Viscoelastic Drops: the Oldroyd-B Model}
\label{sec:oldroydB}
%% ════════════════════════════════════════════════════════════════════════════

Every result in Sections 1--9 assumes the viscous stress is instantaneously proportional to the rate of strain,
$\boldsymbol{\tau} = 2\mu\,\mathbf{e}$. Real polymer solutions violate this.
A drop of polyethylene oxide (PEO) in water, a silicone oil of moderate
molecular weight, or a biological fluid such as mucus all retain a
\emph{memory} of past deformations. Stress does not vanish the instant
deformation stops; it relaxes on a characteristic timescale set by the
polymer microstructure.

The \textbf{Oldroyd-B model} is the simplest constitutive law that captures
this memory and remains analytically tractable. It is the
canonical model for dilute polymer solutions: a Newtonian solvent (viscosity
$\mu_s$) in which elastic polymer dumbbells are suspended (contributing
additional viscosity $\mu_p$, with $\mu = \mu_s + \mu_p$ the total). The
model has two timescales: a \emph{relaxation time} $\lambda_1$ (how long
polymer stress persists after deformation stops) and a \emph{retardation
time} $\lambda_2 < \lambda_1$ (how long the combined fluid takes to respond
to a new deformation). Setting $\lambda_2 = 0$ reduces Oldroyd-B to the
Maxwell model; setting $\lambda_1 = \lambda_2$ recovers a Newtonian fluid.

Reid's characteristic equation \eqref{eq:char_eq} survives generalisation
to Oldroyd-B. The left-hand side $\alpha^4/q^4 + 1$ is unchanged. The
right-hand side is Reid's right-hand side with $q$ replaced by a modified
wavenumber $q_*$, a rational function of $q^2$. This rational structure
introduces a new type of pole, the \emph{retardation pole}, which captures
additional eigenvalues beyond Reid's fundamental pair and eventually causes
the two-number $\{A_l, D_l\}$ picture to break down.

%% ────────────────────────────────────────────────────────────────────────────
\subsection{The Oldroyd-B constitutive law}
%% ────────────────────────────────────────────────────────────────────────────

\subsubsection{Physical motivation: the mechanical analog}

\begin{notebox}[Spring-dashpot picture of Oldroyd-B]
Imagine the fluid as two \emph{parallel} stress-carrying branches:
\begin{itemize}
  \item \textbf{Solvent branch:} a pure dashpot of viscosity $\mu_s$.
        Stress responds instantaneously to strain rate: $\boldsymbol{\tau}_s
        = 2\mu_s\,\mathbf{e}$. No memory.
  \item \textbf{Polymer branch:} a Hookean spring (elastic modulus
        $G = \mu_p/\lambda_1$) in \emph{series} with a dashpot of
        viscosity $\mu_p$. The spring stores elastic energy when the
        dumbbell is stretched; the dashpot relaxes that energy on timescale
        $\lambda_1 = \mu_p/G$.
\end{itemize}
The total stress is $\boldsymbol{\tau} = \boldsymbol{\tau}_s +
\boldsymbol{\tau}_p$. The retardation time
\begin{equation}
  \lambda_2 = \lambda_1\,\frac{\mu_s}{\mu_s + \mu_p} = \beta_s\,\lambda_1
\end{equation}
is the timescale on which the \emph{combined} system (solvent + polymer)
first responds to a new deformation: the solvent wants to respond
immediately, but the spring in the polymer arm initially resists,
delaying the total response.
\end{notebox}

\subsubsection{The constitutive equation: two equivalent forms}

In terms of the two raw timescales $(\lambda_1, \lambda_2)$:
\begin{equation}
  \boldsymbol{\tau} + \lambda_1\,\overset{\nabla}{\boldsymbol{\tau}}
  = 2\mu\!\left(\mathbf{e} + \lambda_2\,\overset{\nabla}{\mathbf{e}}\right),
  \qquad 0 \leq \lambda_2 < \lambda_1,
  \label{eq:OB_raw}
\end{equation}
where $\overset{\nabla}{(\cdot)}$ is the upper-convected time derivative:
\begin{equation}
  \overset{\nabla}{\mathbf{A}}
  = \frac{\partial\mathbf{A}}{\partial t}
  + (\bu\cdot\nabla)\mathbf{A}
  - (\nabla\bu)\cdot\mathbf{A}
  - \mathbf{A}\cdot(\nabla\bu)^\top.
  \label{eq:ucd}
\end{equation}
In terms of the retardation ratio $\beta_s \equiv \lambda_2/\lambda_1$:
\begin{equation}
  \boldsymbol{\tau} + \lambda_1\,\overset{\nabla}{\boldsymbol{\tau}}
  = 2\mu\!\left(\mathbf{e} + \beta_s\lambda_1\,\overset{\nabla}{\mathbf{e}}\right),
  \qquad \beta_s \in [0,1),
  \label{eq:OB_beta}
\end{equation}
with $\beta_s = \mu_s/(\mu_s+\mu_p)$ the solvent viscosity fraction.
The two forms \eqref{eq:OB_raw} and \eqref{eq:OB_beta} are identical;
we carry both throughout because each makes different limits transparent.

\begin{notebox}[Upper-convected derivative at linearised order]
All fields are $O(\epsilon)$. In the upper-convected derivative
\eqref{eq:ucd}, the terms $(\bu\cdot\nabla)\mathbf{A}$,
$(\nabla\bu)\cdot\mathbf{A}$, and $\mathbf{A}\cdot(\nabla\bu)^\top$ are
each products of two $O(\epsilon)$ quantities, hence $O(\epsilon^2)$.
At linear order:
\[
  \overset{\nabla}{\boldsymbol{\tau}} \;\longrightarrow\;
  \frac{\partial\boldsymbol{\tau}}{\partial t} + O(\epsilon^2), \qquad
  \overset{\nabla}{\mathbf{e}} \;\longrightarrow\;
  \frac{\partial\mathbf{e}}{\partial t} + O(\epsilon^2).
\]
We record the full upper-convected form for completeness and frame-
indifference, but drop the $O(\epsilon^2)$ terms immediately.
\end{notebox}

\subsubsection{Parameter limits}

\begin{center}
\begin{tabular}{lllll}
  \toprule
  Limit & Condition & Model & $\beta_s$ & \\
  \midrule
  Maxwell    & $\lambda_2 = 0$        & Single-relaxation memory & $0$     & \\
  Newtonian  & $\lambda_1 = \lambda_2$& Viscous, viscosity $\mu$ & $1$     & \\
  Solvent    & $\lambda_1\to 0$ at fixed $G=\mu_p/\lambda_1$, $\mu_s$
                                      & Newtonian, viscosity $\mu_s$ & $\to 1$ & \\
  Elastic solid & $\lambda_1\to\infty$, $\beta_s=0$, $G=\mu_p/\lambda_1$ fixed ($\mu_p\to\infty$)
                                      & Hookean solid, modulus $G$ & $0$ & \\
  \bottomrule
\end{tabular}
\end{center}

%% ────────────────────────────────────────────────────────────────────────────
\subsection{Linearisation and effective viscosity}
%% ────────────────────────────────────────────────────────────────────────────

The linearisation strategy in this section follows Zrni\'{c} \& Brenn
\cite{zirnic2024}, whose Eq.\,(3.3) establishes the same effective-viscosity
replacement $\mu \to \mu_{\mathrm{eff}}(\sigma)$ for the free-surface Oldroyd-B
problem.

\begin{derivbox}[Linearised Oldroyd-B in the frequency domain]
With time dependence $e^{-\sigma t}$ (our convention: $\mathrm{Re}(\sigma)>0$
for decay), $\partial_t \to -\sigma$. Dropping $O(\epsilon^2)$ upper-convected
terms:
\[
  (1 - \lambda_1\sigma)\,\boldsymbol{\tau}
  = 2\mu\,(1 - \lambda_2\sigma)\,\mathbf{e}.
\]
Dividing through by $(1 - \lambda_1\sigma)$ (valid generically, since
$\sigma = 1/\lambda_1$ is a special value examined separately):
\begin{equation}
  \boldsymbol{\tau} = 2\mu_{\mathrm{eff}}(\sigma)\,\mathbf{e},
  \qquad
  \mu_{\mathrm{eff}}(\sigma) = \mu\,\frac{1-\lambda_2\sigma}{1-\lambda_1\sigma}.
  \label{eq:mueff_OB}
\end{equation}
In both parameterisations simultaneously:
\begin{equation}
  \mu_{\mathrm{eff}}(\sigma)
  = \mu\,\frac{1-\lambda_2\sigma}{1-\lambda_1\sigma}
  = \mu\,\frac{1-\beta_s\lambda_1\sigma}{1-\lambda_1\sigma}.
  \label{eq:mueff_both}
\end{equation}
The effective kinematic viscosity is $\nu_{\mathrm{eff}} =
\mu_{\mathrm{eff}}/\rho$.
\end{derivbox}

\begin{notebox}[Analytic structure of $\mu_{\mathrm{eff}}(\sigma)$]
As a function of $\sigma$, $\mu_{\mathrm{eff}}$ is a ratio of two linear
factors:
\begin{itemize}
  \item \textbf{Zero} at $\sigma = 1/\lambda_2$: solvent and polymer stresses
        cancel in the tangential direction. ($\lambda_2 = 0$ pushes this
        zero to infinity (the Maxwell model has no zero.)
  \item \textbf{Pole} at $\sigma = 1/\lambda_1$: polymer stress cannot relax;
        effective viscosity diverges (solid-like response).
\end{itemize}
Setting $\lambda_2 = 0$: $\mu_{\mathrm{eff}} = \mu/(1-\lambda_1\sigma)$,
recovering the Maxwell effective viscosity. Setting $\lambda_1 = \lambda_2$:
$\mu_{\mathrm{eff}} = \mu$, recovering Newtonian. The Oldroyd-B model
interpolates between these two limits.
\end{notebox}

\subsubsection{The two Deborah numbers}

Referenced to the inviscid oscillation frequency $\sigma_{l;0}$:
\begin{equation}
  De_1 \equiv \lambda_1\,\sigma_{l;0},
  \qquad
  De_2 \equiv \lambda_2\,\sigma_{l;0} = \beta_s\,De_1.
  \label{eq:De_defs}
\end{equation}
$De_j$ is the ratio of the fluid's $j$th memory timescale to the drop's
natural oscillation period.

\begin{notebox}[Convention note]
Zrni\'{c} \& Brenn \cite{zirnic2024} define their Deborah numbers referenced
to the capillary timescale $(T_1/(\rho R^3))^{1/2}$, whereas here they are
referenced to $\sigma_{l;0} = [l(l-1)(l+2)]^{1/2}(T_1/(\rho R^3))^{1/2}$.
Consequently $De_j^{\text{here}} = \sqrt{l(l-1)(l+2)}\,De_j^{\text{Z\&B}}$;
numerical values are mode-dependent and should not be compared directly.
\end{notebox} When $De_1 \ll 1$, the polymer relaxes many times
per cycle and the drop behaves as Newtonian. When $De_1 \sim 1$, elasticity
and surface tension compete on equal footing.

\begin{derivbox}[Expressing $\lambda_j\sigma$ in terms of $q^2$ and $\alpha^2$]
Recall $\sigma = q^2\nu/R^2$ and $\sigma_{l;0} = \alpha^2\nu/R^2$.
Therefore:
\begin{equation}
  \lambda_j\sigma
  = \lambda_j\cdot\frac{q^2\nu}{R^2}
  = \frac{\lambda_j\sigma_{l;0}}{\alpha^2\nu/R^2}\cdot\frac{q^2\nu}{R^2}
  = \frac{De_j\cdot q^2}{\alpha^2}.
  \label{eq:lamsig}
\end{equation}
Substituting into \eqref{eq:mueff_OB}:
\begin{equation}
  \mu_{\mathrm{eff}} = \mu\,\frac{\alpha^2 - De_2\,q^2}{\alpha^2 - De_1\,q^2},
  \qquad
  \nu_{\mathrm{eff}} = \nu\,\frac{\alpha^2 - De_2\,q^2}{\alpha^2 - De_1\,q^2}.
  \label{eq:nueff_dimless}
\end{equation}
All dimensional quantities have cancelled. The effective viscosity is a
function of $q^2/\alpha^2$ only.
\end{derivbox}

%% ────────────────────────────────────────────────────────────────────────────
\subsection{Boundary condition audit}
%% ────────────────────────────────────────────────────────────────────────────

We have a free-surface drop in air. We examine each of the three boundary
conditions.

\subsubsection{BC1 (kinematic): unchanged}

The kinematic condition equates the fluid radial velocity at $x=1$ to the
rate of change of the surface position. It is a purely geometric statement
involving only the \emph{velocity field}, not the stress. The Oldroyd-B
constitutive law does not alter the velocity field structure. Therefore:
\begin{equation}
  U(1) = -1. \tag{BC1, unchanged}
\end{equation}

\subsubsection{BC2 (tangential stress): unchanged in form}

The free-surface condition $\tau_{r\theta}|_{x=1} = 0$ becomes:
\[
  \mu_{\mathrm{eff}}(\sigma)\!\left[\,r\,\frac{\partial}{\partial r}
  \!\left(\frac{u_\theta}{r}\right)
  + \frac{1}{r}\frac{\partial u_r}{\partial\theta}\,\right]_{r=R} = 0.
\]
Since $\mu_{\mathrm{eff}} \neq 0$ generically, we divide through. The
condition reduces to exactly the same operator equation on $U$:
\begin{equation}
  \left[\,\frac{\dd^2}{\dd x^2}
  - \frac{2}{x}\frac{\dd}{\dd x}
  + \frac{l(l+1)}{x^2}\,\right]U\,\bigg|_{x=1} = 0. \tag{BC2, unchanged}
\end{equation}

\begin{notebox}[The zero of $\mu_{\mathrm{eff}}$ at $\sigma = 1/\lambda_2$]
At $\sigma = 1/\lambda_2$, the effective viscosity $\mu_{\mathrm{eff}} \to 0$
and BC2 becomes $0 = 0$ for any $U$: it imposes no constraint.
The frequency $\sigma = 1/\lambda_2$ is the \emph{retardation pole} of
$q_*^2$: as $w \to w_{\mathrm{ret}} = \alpha^2/De_2$, $q_*^2 \to \infty$.
It is not itself a root of the characteristic equation $F(w) = 0$. The actual retardation roots instead lie in a neighbourhood of $w_{\mathrm{ret}}$ (see Section~\ref{sec:OB_roots}).
At those roots BC2 must be handled by a limiting argument. In practice these
roots decay at rate $1/\lambda_2$, which is faster than the capillary
timescale when $De_2 < 1$, and they are discarded as overtones. For
$De_2 > 1$ a more careful treatment is required.
\end{notebox}

\subsubsection{BC3 (normal stress): viscoelastic correction check}

\begin{derivbox}[First normal stress difference at $O(\epsilon)$ for Oldroyd-B]
For the Oldroyd-B family, the first normal stress difference is:
\[
  N_1 = \tau_{rr} - \tau_{\theta\theta}.
\]
At the linearised level, $\boldsymbol{\tau} = 2\mu_{\mathrm{eff}}\,\mathbf{e}$
with $\mathbf{e} = O(\epsilon)$, so $N_1 = O(\epsilon)$. Does this $O(\epsilon)$
quantity add a new term to BC3?

The normal stress balance at a free surface equates the jump in
\emph{radial} normal stress to the Laplace pressure:
\[
  -p_{rr}\big|_{\text{inside}} = T_1\!\left(\frac{1}{R_1}+\frac{1}{R_2}\right).
\]
The radial normal stress inside is $-p_{rr} = p + \delta p - \tau_{rr}$,
with $\tau_{rr} = 2\mu_{\mathrm{eff}}\,\partial u_r/\partial r$. This is
precisely the viscous term already present in BC3 \eqref{eq:BC3}, now with
$\mu \to \mu_{\mathrm{eff}}$. The remaining part $N_1 = \tau_{rr} -
\tau_{\theta\theta}$ measures the \emph{anisotropy} of the stress tensor,
but it does not enter BC3 at $O(\epsilon)$.
The reason is that the base state is a quiescent sphere with
$\boldsymbol{\tau}_0 = 0$, so the $O(\epsilon)$ tilt of the outward normal
from $\hat{\mathbf{r}}$ introduces only an $O(\epsilon)$ correction to the
normal direction times an $O(\epsilon^0) = 0$ base tangential traction,
contributing nothing at $O(\epsilon)$.
Therefore $N_1$ contributes nothing extra to BC3 at $O(\epsilon)$.

At $O(\epsilon^2)$, the surface is deformed and the mismatch between normal
and tangential directions at the perturbed surface generates an
$O(\epsilon^2)$ correction from $N_1$ to the effective surface tension.
This is an open problem noted in Section~\ref{sec:OB_open}.
\end{derivbox}

The normal stress boundary condition is:
\begin{equation}
  (l-1)(l+2)\frac{T_1}{\rho R}
  = P_0 - 2\nu_{\mathrm{eff}}(\sigma)\,\sigma\!\left[\,U'(1) - 2U(1)\,\right].
  \tag{BC3$'$}
  \label{eq:BC3_OB}
\end{equation}
This is identical to the Newtonian BC3 \eqref{eq:BC3} with
$\nu \to \nu_{\mathrm{eff}}$ from \eqref{eq:nueff_dimless}.

%% ────────────────────────────────────────────────────────────────────────────
\subsection{Modified governing ODE}
%% ────────────────────────────────────────────────────────────────────────────

\subsubsection{The modified wavenumber $q_*$}

\begin{derivbox}[Definition and derivation of $q_*^2$]
In Reid's Newtonian derivation, the parameter $q^2 = \sigma R^2/\nu$ emerged
when the $r$-momentum equation was multiplied through by $R/(\nu\sigma)$.
For Oldroyd-B, the same step with $\nu \to \nu_{\mathrm{eff}}$ gives:
\begin{equation}
  q_*^2 \;\equiv\; \frac{\sigma R^2}{\nu_{\mathrm{eff}}(\sigma)}
  = q^2\,\frac{1-\lambda_1\sigma}{1-\lambda_2\sigma}.
  \label{eq:qstar_raw}
\end{equation}
Using \eqref{eq:lamsig} to express $\lambda_j\sigma$ in terms of $q^2$
and $\alpha^2$:
\begin{equation}
  \boxed{
    q_*^2 = q^2\cdot\frac{\alpha^2 - De_1\,q^2}{\alpha^2 - De_2\,q^2}.
  }
  \label{eq:qstar_OB}
\end{equation}
In both parameterisations:
\begin{equation}
  q_*^2 = q^2\cdot\frac{\alpha^2 - De_1\,q^2}{\alpha^2 - \beta_s De_1\,q^2}.
  \label{eq:qstar_both}
\end{equation}
\end{derivbox}

\begin{notebox}[Structural difference from the Maxwell model]
For the Maxwell fluid, $q_*^2 = q^2(1 - De_1 q^2/\alpha^2) =
q^2(\alpha^2 - De_1 q^2)/\alpha^2$: the denominator is the constant
$\alpha^2$. For Oldroyd-B, the denominator is $\alpha^2 - De_2 q^2$,
which \emph{vanishes} at $q^2 = \alpha^2/De_2$, i.e.\ at
$\sigma = \sigma_{l;0}/De_2 = 1/\lambda_2$. As $q^2 \to \alpha^2/De_2$:
$q_*^2 \to \infty$. This is the \textbf{retardation pole}: a singularity
in the effective wavenumber absent from Maxwell or Newtonian problems. All qualitative differences between Oldroyd-B and Maxwell originate in this pole.
\end{notebox}

\begin{notebox}[Limits of $q_*^2$]
\begin{itemize}
  \item $De_2 = 0$ (Maxwell): $q_*^2 = q^2(\alpha^2 - De_1 q^2)/\alpha^2$.
        Retardation pole disappears.
  \item $De_1 = De_2$ (Newtonian): numerator equals denominator,
        $q_*^2 = q^2$.
  \item $q^2 = \alpha^2/De_1$ (relaxation zero): numerator $= 0$,
        $q_*^2 = 0$. At $\sigma = 1/\lambda_1$ the polymer is fully
        relaxed; it contributes zero stress, as if absent.
  \item $q^2 = \alpha^2/De_2$ (retardation pole): denominator $= 0$,
        $q_*^2\to\infty$. At $\sigma = 1/\lambda_2$ the effective
        viscosity vanishes and the Bessel argument diverges.
\end{itemize}
Since $De_2 < De_1$, the retardation pole $\alpha^2/De_2 > \alpha^2/De_1$
lies to the right of the relaxation zero on the positive real $q^2$-axis.
\end{notebox}

\subsubsection{The ODE for $U(x)$ under Oldroyd-B}

\begin{derivbox}[The pressure field is unchanged]
The pressure perturbation satisfies $\nabla^2(\delta p/\rho) = 0$ regardless
of $\nu_{\mathrm{eff}}$. To see this: take the divergence of the linearised
momentum equation
$-\sigma\bu = -\nabla(\delta p/\rho) + \nu_{\mathrm{eff}}\nabla^2\bu$
and use incompressibility $\nabla\cdot\bu = 0$. Since $\nu_{\mathrm{eff}}$
is a scalar constant (it depends on $\sigma$ but not on $\mathbf{x}$), we get
$\nabla^2(\delta p/\rho) = \nu_{\mathrm{eff}}\nabla^2(\nabla\cdot\bu) = 0$.
The solution $\delta p/\rho = \epsilon P_0 x^l Y_l^m$ is unchanged.
\end{derivbox}

Following Section~5 with the substitution $\nu \to \nu_{\mathrm{eff}}$
and hence $q^2 \to q_*^2$:

\begin{keybox}[ODE for $U(x)$ under Oldroyd-B]
\begin{equation}
  \left[\,\frac{\dd^2}{\dd x^2} - \frac{l(l+1)}{x^2} + q_*^2\,\right]U(x)
  = q_*^2\,\Pi_0\,x^{l+1},
  \label{eq:ODE_OB}
\end{equation}
with $q_*^2$ given by \eqref{eq:qstar_OB}. The general solution is:
\begin{equation}
  U(x) = C\,x\,j_l(q_* x) + \Pi_0\,x^{l+1}.
  \label{eq:U_OB}
\end{equation}
This is Reid's ODE and general solution \eqref{eq:ODE_U}--\eqref{eq:U_general}
with the single replacement $q \to q_*$.
\end{keybox}

%% ────────────────────────────────────────────────────────────────────────────
\subsection{The Oldroyd-B characteristic equation}
%% ────────────────────────────────────────────────────────────────────────────

\subsubsection{The ratio cancellation}

Before deriving the characteristic equation, we establish a key algebraic
fact that determines its structure.

\begin{derivbox}[The ratio $\alpha_*^2/q_*^2$ is Newtonian]
Define, in analogy with \eqref{eq:alpha_def}:
\begin{equation}
  \alpha_*^2 \;\equiv\; \frac{\sigma_{l;0}R^2}{\nu_{\mathrm{eff}}}
  = \alpha^2\,\frac{1-\lambda_1\sigma}{1-\lambda_2\sigma}
  = \alpha^2\cdot\frac{\alpha^2-De_1\,q^2}{\alpha^2-De_2\,q^2}.
  \label{eq:alphastar_OB}
\end{equation}
Computing the ratio $\alpha_*^2/q_*^2$:
\[
  \frac{\alpha_*^2}{q_*^2}
  = \frac{\displaystyle\alpha^2\cdot\frac{\alpha^2-De_1 q^2}{\alpha^2-De_2 q^2}}
         {\displaystyle q^2\cdot\frac{\alpha^2-De_1 q^2}{\alpha^2-De_2 q^2}}
  = \frac{\alpha^2}{q^2}.
\]
The Oldroyd-B factors cancel exactly in the ratio. Therefore:
\begin{equation}
  \frac{\alpha_*^4}{q_*^4} = \frac{\alpha^4}{q^4}.
  \label{eq:ratio_cancel}
\end{equation}
\end{derivbox}

\begin{notebox}[Why the cancellation occurs]
The ratio $\alpha^4/q^4 = \sigma_{l;0}^2/\sigma^2$ is a pure ratio of
frequencies. It is independent of viscosity by construction, since both
$\alpha^2 \propto 1/\nu$ and $q^2 \propto 1/\nu$ rescale identically when
$\nu \to \nu_{\mathrm{eff}}$. The cancellation \eqref{eq:ratio_cancel} holds
for \emph{any} constitutive model that linearises to $\boldsymbol{\tau} =
2\mu_{\mathrm{eff}}(\sigma)\,\mathbf{e}$ with a scalar $\mu_{\mathrm{eff}}$.
It is a universal feature of the linearised problem, not specific to
Oldroyd-B.
\end{notebox}

\subsubsection{Eliminating $T_1$}

\begin{derivbox}[Modified rescaling]
Multiplying BC3$'$ \eqref{eq:BC3_OB} through by $lR^2/\nu_{\mathrm{eff}}^2$
and using $(l-1)(l+2)T_1/(\rho R) = \sigma_{l;0}^2 R^2/l$:
\[
  \underbrace{\frac{\sigma_{l;0}^2 R^4}{\nu_{\mathrm{eff}}^2}}_{\alpha_*^4}
  = \underbrace{\frac{\sigma^2 R^4}{\nu_{\mathrm{eff}}^2}}_{q_*^4}\Pi_0
  - 2l\underbrace{\frac{\sigma R^2}{\nu_{\mathrm{eff}}}}_{q_*^2}\!\left[U'(1)-2U(1)\right].
\]
Hence:
\begin{equation}
  \alpha_*^4 = q_*^2\!\left\{\,q_*^2\Pi_0
  - 2l\!\left[U'(1) - 2U(1)\right]\right\}.
  \label{eq:alpha4star_OB}
\end{equation}
$T_1$ and $\rho$ have cancelled completely. By \eqref{eq:ratio_cancel},
$\alpha_*^4 = q_*^4\cdot\alpha^4/q^4$, so the left-hand side of the
characteristic equation will be $\alpha^4/q^4$, identical to the Newtonian
case.
\end{derivbox}

\subsubsection{Applying BC1 and BC2 and assembling the characteristic equation}

The algebra of Section~7.2 carries through verbatim with $q \to q_*$:

\paragraph{BC1:} $U(1) = C\,j_l(q_*) + \Pi_0 = -1.$

\paragraph{BC2:} Evaluating \eqref{eq:BC2} on \eqref{eq:U_OB}:
\[
  C\!\left[(2l-q_*^2)\,j_l(q_*) - 2q_*\,j_l'(q_*) + 2(l^2-1)\,j_l(q_*)\right]
  + 2(l^2-1)\Pi_0 = 0.
\]

\paragraph{Solving for $C$ and $\Pi_0$:}
Substituting $\Pi_0 = -1 - C j_l(q_*)$:
\begin{equation}
  C = \frac{2(l-1)(l+1)}{(2l-q_*^2)\,j_l(q_*) - 2q_*\,j_l'(q_*)}.
  \label{eq:C_OB}
\end{equation}
Using the Bessel derivative identity $q_* j_l'(q_*)/j_l(q_*) =
l - q_* Q_{l+1/2}(q_*)$ and the Bessel recurrence to simplify, then
substituting $C$, $\Pi_0$, $U'(1)$, and $U(1)$ into \eqref{eq:alpha4star_OB}:

\begin{keybox}[Oldroyd-B Characteristic Equation]
\begin{equation}
  \frac{\alpha^4}{q^4} + 1
  = \frac{2(l-1)}{q_*^2}
    \left[\,l + (l+1)\,\frac{q_* - 2l\,Q_{l+1/2}(q_*)}{q_* - 2\,Q_{l+1/2}(q_*)}\,\right],
  \label{eq:char_OB}
\end{equation}
where
\begin{equation}
  q_*^2 = q^2\cdot\frac{\alpha^2-De_1\,q^2}{\alpha^2-De_2\,q^2}
  = q^2\cdot\frac{\alpha^2-De_1\,q^2}{\alpha^2-\beta_s De_1\,q^2},
  \label{eq:qstar_OB_final}
\end{equation}
$Q_{l+1/2}(q_*) = J_{l+3/2}(q_*)/J_{l+1/2}(q_*)$, and the parameters
$\alpha^2 = \sigma_{l;0}R^2/\nu$, $De_1 = \lambda_1\sigma_{l;0}$,
$De_2 = \lambda_2\sigma_{l;0} = \beta_s De_1$ are real constants.

\smallskip\noindent Limits:
$De_1 = 0$ (or $\beta_s = 1$): $q_* = q$, recovers Reid \eqref{eq:char_eq}.\\
$De_2 = 0$: recovers the Maxwell characteristic equation.\\
$De_1 = De_2 \neq 0$: $q_* = q$, again recovers Reid \eqref{eq:char_eq}.\\[2pt]
This equation coincides with Zrni\'{c} \& Brenn \cite{zirnic2024} Eqs.\,(3.11)--(3.12).
\end{keybox}

\begin{notebox}[What changed from Reid's equation]
\begin{enumerate}
  \item The left-hand side $\alpha^4/q^4 + 1$ is \emph{identical} to Reid's.
        (Consequence of the ratio cancellation \eqref{eq:ratio_cancel}.)
  \item The right-hand side is Reid's right-hand side with $q \to q_*$,
        where $q_*$ is the rational function \eqref{eq:qstar_OB_final}.
  \item The entire Oldroyd-B physics is encoded in $De_1$ and $\beta_s$
        (two extra real parameters beyond Reid's $\alpha^2$).
  \item The rational structure of $q_*^2$ introduces a new pole
        at $q^2 = \alpha^2/De_2$ with no Newtonian or Maxwell analogue.
\end{enumerate}
\end{notebox}

%% ────────────────────────────────────────────────────────────────────────────
\subsection{Structure of solutions}
\label{sec:OB_roots}
%% ────────────────────────────────────────────────────────────────────────────

\subsubsection{The characteristic equation as a condition in $q^2$}

Let $w = q^2$. Define the rational map:
\begin{equation}
  R(w) \;\equiv\; q_*^2
  = w\cdot\frac{\alpha^2 - De_1 w}{\alpha^2 - De_2 w}.
  \label{eq:R_map}
\end{equation}
Then \eqref{eq:char_OB} is a transcendental equation $F(w) = 0$ in the single
complex variable $w$, with $\alpha^2$, $De_1$, $De_2$ as real parameters.
Its roots $w_j$ give eigenvalues $\sigma_j = w_j\nu/R^2$.

\subsubsection{Two types of poles in $F(w)$}
\label{sec:two_pole_types}

The function $F(w)$ has poles wherever the right-hand side of
\eqref{eq:char_OB} diverges. There are two distinct sources:

\paragraph{Type 1: Bessel poles.}
The right-hand side of \eqref{eq:char_OB} has poles wherever
$J_{l+1/2}(q_*) = 0$, i.e.\ at the zeros $z_1^* < z_2^* < \cdots$ of the
Bessel function in the $q_*$-plane. These satisfy $q_* = z_n^*$, i.e.\
$R(w) = (z_n^*)^2$: for each $n$, this is a quadratic (in $w$) algebraic
equation with up to two solutions in the $w$-plane. For small $De_1$ the rational map $R(w)$ is a small perturbation of the
Newtonian case, and a heuristic application of the argument principle suggests
that $F(w) = 0$ retains exactly two roots per strip, as in the Newtonian
problem; a rigorous contour argument for general $De_1$ is open (see
\S\ref{sec:OB_open}). The first strip ($n=1$) captures the
\textbf{dominant pair}, governing the fundamental capillary oscillation.

\paragraph{Type 2: Retardation pole.}
The factor $\alpha^2 - De_2 w$ in the denominator of $R(w)$ vanishes at:
\begin{equation}
  w_{\mathrm{ret}} = \frac{\alpha^2}{De_2},
  \qquad
  \sigma_{\mathrm{ret}} = \frac{\sigma_{l;0}}{De_2} = \frac{1}{\lambda_2}.
  \label{eq:ret_pole}
\end{equation}
As $w \to w_{\mathrm{ret}}$, $q_* \to \infty$ and $Q_{l+1/2}(q_*)$ oscillates
rapidly. Near this singularity the characteristic equation
\eqref{eq:char_OB} can be satisfied, capturing one or two
\textbf{retardation roots}. These correspond to modes decaying at rate
$\sim 1/\lambda_2$. Since $\lambda_2 < \lambda_1$, the retardation mode
always decays faster than the Maxwell-type relaxation mode; whether it is
faster or slower than the fundamental capillary pair depends on $De_2$.

\begin{notebox}[Physical meaning of the retardation roots]
At $\sigma = 1/\lambda_2$, the solvent and polymer contributions to the
tangential stress cancel: $\mu_{\mathrm{eff}}(1/\lambda_2) = 0$. The
retardation roots describe a mode in which the polymer network, instead of
resisting the flow, briefly \emph{drives} it, before the solvent dissipates
the energy. This phenomenon has no Newtonian analogue; it requires two fluid components. Physically, it corresponds to the elastic recoil
of the polymer network overshooting the equilibrium and being damped by the
solvent on the retardation timescale $\lambda_2$.
\end{notebox}

\subsubsection{Dominance condition for the fundamental pair}

The fundamental pair of roots governs observable drop oscillation provided
the retardation roots are faster-decaying (subdominant):
\begin{equation}
  \mathrm{Re}(\sigma_{1,2}) < \frac{1}{\lambda_2}
  \quad\Longleftrightarrow\quad
  De_2\cdot\mathrm{Re}(\beta_{1,2}) < 1,
  \label{eq:dominance}
\end{equation}
where $\beta_{j} = q_j^2/\alpha^2 = \sigma_j/\sigma_{l;0}$.

When \eqref{eq:dominance} holds, the retardation roots are higher overtones:
they decay faster than the capillary oscillation period and are irrelevant to
the Mola\v{c}ek \& Bush quasi-static model. When \eqref{eq:dominance} fails,
the retardation roots must be retained.

%% ────────────────────────────────────────────────────────────────────────────
\subsection{$A_l$ and $D_l$ for the Oldroyd-B fluid}
%% ────────────────────────────────────────────────────────────────────────────

\subsubsection{Vieta extraction: the two-number picture survives}

Provided \eqref{eq:dominance} holds, the dynamics of mode $l$ is governed by
the two roots $\beta_{1,2}$ of the characteristic equation \eqref{eq:char_OB}
in the first Bessel strip. The Vieta extraction of Section~9 applies without
modification:
\begin{equation}
  A_l^{\mathrm{OB}} = \frac{1}{\beta_1\beta_2},
  \qquad
  D_l^{\mathrm{OB}} = \frac{\beta_1 + \beta_2}{2a\,\beta_1\beta_2},
  \qquad
  a = \mathrm{Oh}\sqrt{l(l-1)(l+2)}.
  \label{eq:AD_OB}
\end{equation}
The difference from the Newtonian case is that $\beta_{1,2}$ are now roots
of \eqref{eq:char_OB} rather than \eqref{eq:char_eq}: they depend on
$(De_1, \beta_s)$ in addition to $(l, \mathrm{Oh})$.

\begin{keybox}[Parameter dependence of $A_l^{\mathrm{OB}}$ and $D_l^{\mathrm{OB}}$]
\begin{equation}
  A_l^{\mathrm{OB}} = A_l^{\mathrm{OB}}(l,\,\mathrm{Oh},\,De_1,\,\beta_s),
  \qquad
  D_l^{\mathrm{OB}} = D_l^{\mathrm{OB}}(l,\,\mathrm{Oh},\,De_1,\,\beta_s).
\end{equation}
Limits:
\begin{itemize}
  \item $De_1 \to 0$ (or $\beta_s \to 1$): recover Mola\v{c}ek \& Bush
        Newtonian values.
  \item $\beta_s \to 0$ at fixed $De_1$: recover Maxwell values (a special
        case of the Oldroyd-B framework).
  \item $De_1 \to \infty$ at fixed $\beta_s$: strongly elastic limit; both
        coefficients must be determined numerically.
\end{itemize}
\end{keybox}

\begin{notebox}[Monotone inequalities between models (conjectured)]
For the same total viscosity $\mu = \mu_s + \mu_p$ and the same Oh,
numerical evidence suggests:
\[
  A_l^{\mathrm{N}} \;\leq\; A_l^{\mathrm{OB}} \;\leq\; A_l^{\mathrm{M}},
  \qquad
  D_l^{\mathrm{M}} \;\leq\; D_l^{\mathrm{OB}} \;\leq\; D_l^{\mathrm{N}}.
\]
Elasticity enhances effective inertia (increases $A_l$) and reduces
instantaneous dissipation (decreases $D_l$). The solvent, present in
Oldroyd-B but absent in Maxwell, partially restores both toward their
Newtonian values. The interpolation appears monotone in $\beta_s$: as
$\beta_s$ increases from 0 (Maxwell) to 1 (Newtonian), $A_l$ decreases
and $D_l$ increases. These inequalities are consistent with physical
intuition but have not been established analytically from the
transcendental characteristic equation; a formal proof is open.
\end{notebox}

\subsubsection{When the two-number picture breaks}

When condition \eqref{eq:dominance} fails, the retardation roots have decay
rates comparable to the fundamental pair. There are then three (or more) roots
of comparable magnitude, and the Mola\v{c}ek \& Bush quadratic ODE for $b_l$
is no longer adequate. In capillary time $\tau = \sigma_{l;0}\,t$ (all quantities
dimensionless), the full description is an integro-differential equation,
\begin{equation}
  A_l^{\mathrm{OB}}\,\ddot{b}_l
  + 2a
    \int_0^\tau K^{(\mathrm{cap})}(\tau-\tau')\,\dot{b}_l(\tau')\,\dd\tau'
  + b_l = 0, \qquad a \equiv \mathrm{Oh}\sqrt{l(l-1)(l+2)},
  \label{eq:integrodiff_OB}
\end{equation}
where dots denote $\dd/\dd\tau$ and the dimensionless OB memory kernel in
capillary time is
\begin{equation}
  K^{(\mathrm{cap})}(\tau) = \beta_s\,\delta(\tau)
  + \frac{1-\beta_s}{De_1}\,e^{-\tau/De_1}.
  \label{eq:OB_kernel_cap}
\end{equation}
Its positive-exponent Laplace transform is
$\tilde{K}^{(\mathrm{cap})}(\beta)
  = \int_0^\infty K^{(\mathrm{cap})}(\tau)\,e^{\beta\tau}\,\dd\tau
  = \nu_{\mathrm{eff}}(\beta\,\sigma_{l;0})/\nu
  = (1 - De_2\,\beta)/(1 - De_1\,\beta)$,
and its zeroth moment is $\int_0^\infty K^{(\mathrm{cap})}\,\dd\tau = \beta_s + (1-\beta_s) = 1$.
In the Newtonian limit ($De_1\to 0$, or equivalently $\beta_s\to 1$):
$K^{(\mathrm{cap})}\to\delta(\tau)$, and the equation reduces to
$A_l^{\mathrm{OB}}\ddot{b}+2a\dot{b}+b=0$.
This is a structured approximation. The convolution kernel replaces the
full Newtonian damping function $D_l(\mathrm{Oh})$ by the zeroth moment
$\tilde{K}^{(\mathrm{cap})}(0)=1$. It captures the frequency-dependent viscosity
$\nu_{\mathrm{eff}}(\beta\sigma_{l;0})/\nu$ but not the Oh-dependence in
Reid's Bessel-function structure. The exact $D_l(\mathrm{Oh})$ is obtained via
Vieta extraction from \eqref{eq:char_OB} at $De_1=0$.
The dimensional kernel $K_{\mathrm{OB}}(s)$ satisfies
$\tilde{K}_{\mathrm{OB}}(\sigma)
  = \int_0^\infty K_{\mathrm{OB}}(s)\,e^{\sigma s}\,\dd s
  = \nu_{\mathrm{eff}}(\sigma)/\nu$:
\begin{equation}
  K_{\mathrm{OB}}(s) = \beta_s\,\delta(s)
  + \frac{1-\beta_s}{\lambda_1}\,e^{-s/\lambda_1},
  \label{eq:OB_kernel}
\end{equation}
related to the capillary kernel by
$K^{(\mathrm{cap})}(\tau) = \sigma_{l;0}^{-1}\,K_{\mathrm{OB}}(\tau/\sigma_{l;0})$.
The $\delta$ term encodes the instantaneous \textbf{solvent} response with
weight $\beta_s = \mu_s/\mu$; the exponential term encodes the polymer memory
with weight $(1-\beta_s) = \mu_p/\mu$ and relaxation time $\lambda_1$.
The two-number approximation corresponds to replacing $K^{(\mathrm{cap})}$ by
its zeroth moment $= 1$ (treating memory as instantaneous), valid when
$De_2 \cdot \mathrm{Re}(\beta_{1,2}) \ll 1$.

%% ────────────────────────────────────────────────────────────────────────────
\subsection{Four clean limits}
%% ────────────────────────────────────────────────────────────────────────────

\begin{center}
\begin{tabular}{lllll}
  \toprule
  Limit & Condition & $q_*^2$ & Char.\ eq.\ & $A_l,D_l$ \\
  \midrule
  Newtonian
    & $De_1 = 0$
    & $q^2$
    & Reid \eqref{eq:char_eq}
    & Mola\v{c}ek \& Bush \\[4pt]
  Maxwell
    & $\beta_s = 0$
    & $q^2\!\left(1 - \tfrac{De_1\,q^2}{\alpha^2}\right)$
    & Khismatullin \& Nadim
    & \eqref{eq:AD_OB}, $\beta_s=0$ \\[4pt]
  Weak elasticity
    & $De_1\ll 1$
    & $q^2[1-(De_1\!-\!De_2)q^2/\alpha^2]$
    & Perturbation of \eqref{eq:char_eq}
    & Analytic correction \\[4pt]
  No retardation
    & $De_2\ll De_1$
    & $\approx q^2(1-De_1 q^2/\alpha^2)$
    & $\approx$ Maxwell
    & $\approx$ Maxwell values \\
  \bottomrule
\end{tabular}
\end{center}

%% ────────────────────────────────────────────────────────────────────────────
\subsection{Open problems and references}
\label{sec:OB_open}
%% ────────────────────────────────────────────────────────────────────────────

\paragraph{What is established.}
\begin{enumerate}
  \item The characteristic equation \eqref{eq:char_OB}--\eqref{eq:qstar_OB_final}
        is exact at $O(\epsilon)$ for a free-surface Oldroyd-B drop in air.
  \item The ratio cancellation \eqref{eq:ratio_cancel} holds for any
        linearisable constitutive model, making the left-hand side of the
        characteristic equation universal.
  \item The dominance condition \eqref{eq:dominance} gives an explicit
        criterion for when the two-number $\{A_l,D_l\}$ extraction is valid.
  \item The integro-differential ODE \eqref{eq:integrodiff_OB}--\eqref{eq:OB_kernel}
        is the natural viscoelastic extension of the Mola\v{c}ek \& Bush quadratic,
        capturing the frequency-dependent damping $\nu_{\mathrm{eff}}/\nu$ when
        the two-number picture breaks down.
\end{enumerate}

\paragraph{What is open.}
\begin{enumerate}
  \item \textbf{Numerical tabulation.} A complete table of
        $A_l^{\mathrm{OB}}(l,\mathrm{Oh},De_1,\beta_s)$ and
        $D_l^{\mathrm{OB}}(l,\mathrm{Oh},De_1,\beta_s)$ analogous to
        Mola\v{c}ek \& Bush Table~1 does not exist in the literature.
        The Newton's method strategy of Section~8 applies directly: for each
        $(De_1,\beta_s)$, initialise from the Newtonian roots and track them
        as $(De_1,\beta_s)$ are increased.
  \item \textbf{Large $De_1$ root structure.} For $De_1 \gg 1$, the numerator
        $\alpha^2 - De_1 q^2$ changes sign for real $q^2 > \alpha^2/De_1$,
        making $q_*^2$ negative there. The physical consequence is that
        Bessel functions $j_l(q_* x)$ become modified spherical Bessel
        functions $i_l(|q_*|x)$, changing the radial profile from
        oscillatory to monotone. A complete argument-principle analysis for
        large $De_1$ does not exist.
  \item \textbf{Second-order normal stress correction.} The $O(\epsilon^2)$
        contribution from $N_1 = \tau_{rr} - \tau_{\theta\theta}$ modifies
        the effective surface tension as
        $T_1 \to T_1^{\mathrm{eff}} = T_1\!\left[1 + \tilde{c}\,\epsilon\,\mathrm{Re}(N_1^{(2)}R/T_1)\right]$
        with $\tilde{c}$ a dimensionless $O(1)$ coefficient (the precise value
        requires the second-order calculation).
        For $\epsilon = O(0.1\text{--}0.3)$ (typical bouncing-drop
        experiments), this correction could reach several percent in $\sigma_{l;0}$.
  \item \textbf{Experimental validation.} No
        experiment has directly measured $A_l$ and $D_l$ for a viscoelastic
        drop and compared to the characteristic equation \eqref{eq:char_OB}.
        PEO-water drops in the range $De_1 \in [0.1, 2]$ and $\beta_s \in
        [0.5, 0.9]$ would be the natural first target.
\end{enumerate}

\paragraph{Key references.}
\begin{itemize}
  \item \textbf{Oldroyd (1950)} \cite{oldroyd1950}: original derivation of
        \eqref{eq:OB_raw}; clearest source for the physical meaning of
        $\lambda_1$, $\lambda_2$, and $\beta_s$.
  \item \textbf{Bird, Armstrong \& Hassager (1987)} \cite{bird1987}:
        dumbbell-model derivation of Oldroyd-B; physical meaning of
        $\beta_s = \mu_s/(\mu_s+\mu_p)$.
  \item \textbf{Khismatullin \& Nadim (2001)} \cite{khismatullin2001}:
        shape oscillations of a viscoelastic drop including the full
        Oldroyd/Jeffreys (relaxation + retardation) constitutive law;
        their characteristic equation already covers the Oldroyd-B case.
  \item \textbf{Zrni\'{c} \& Brenn (2024)} \cite{zirnic2024}:
        free-surface Oldroyd-B drop oscillations; their Eqs.\,(3.11)--(3.12)
        are the same characteristic equation and $q_*^2$ derived in \S\ref{sec:oldroydB}.
        The present notes expand the derivation and connect it to the
        Mola\v{c}ek \& Bush two-number framework.
  \item \textbf{Mukherjee \& Sarkar (2010)} \cite{mukherjee2009}:
        Oldroyd-B drop in a Newtonian bath; numerical results for
        oscillation frequency and damping as functions of $De_1$ and
        $\beta_s$, usable as benchmark for the free-surface limit.
\end{itemize}
